\documentclass[UTF8,12pt]{ctexart}
\usepackage[a4paper,margin=24mm]{geometry}
\usepackage{amsmath,amssymb,bm,graphicx,adjustbox}
\newcommand{\fitmath}[1]{\begin{adjustbox}{max width=\linewidth}$\displaystyle #1$\end{adjustbox}}
\setlength{\parindent}{0pt}
\setlength{\parskip}{.7em}
\sloppy
\allowdisplaybreaks
\begin{document}
\section*{2015 年考研数学三 · 参考答案}
\subsection*{原 PDF 第 25 页}

\subsection*{2015 年真题参考答案}

\subsection*{一、选择题}

(1) D。 (2) C。 (3) B。 (4) C。 (5) D。 (6) A。 (7) C。 (8) B。


\subsection*{二、填空题}

(9) \(\displaystyle -\dfrac{\displaystyle 1}{\displaystyle 2}\)。 (10) \(\displaystyle 2\)。 (11) \(\displaystyle -\dfrac{\displaystyle 1}{\displaystyle 3}dx-\dfrac{\displaystyle 2}{\displaystyle 3}dy\)。 (12) \(\displaystyle 2e^x+e^{-2x}\)。 (13) \(\displaystyle 21\)。 (14) \(\displaystyle \dfrac{\displaystyle 1}{\displaystyle 2}\)。


\subsection*{三、解答题}

(15) \(\displaystyle a=-1,\allowbreak \ b=-\dfrac{\displaystyle 1}{\displaystyle 2},\allowbreak \ k=-\dfrac{\displaystyle 1}{\displaystyle 3}\)。


(16) \(\displaystyle \dfrac{\displaystyle \pi}{\displaystyle 4}-\dfrac{\displaystyle 2}{\displaystyle 5}\)。


(17)（I）证明略。（II）\(\displaystyle 30\)。


(18) \(\displaystyle f(x)=\dfrac{\displaystyle 8}{\displaystyle 4-x},\allowbreak \ x\in I\)。


(19) 证明略。


(20)（I）\(\displaystyle a=0\)。（II）\(\displaystyle X=\begin{pmatrix}3&1&-2\\1&1&-1\\2&1&-1\end{pmatrix}\)。


(21)（I）\(\displaystyle a=4,\allowbreak \ b=5\)。（II）\(\displaystyle P=\begin{pmatrix}1&2&-3\\1&1&0\\-1&0&1\end{pmatrix},\allowbreak \ P^{-1}AP=\begin{pmatrix}5&0&0\\0&1&0\\0&0&1\end{pmatrix}\)。


(22)（I）\(\displaystyle P\{Y=k\}=\dfrac{\displaystyle 1}{\displaystyle 64}(k-1)\left(\dfrac{\displaystyle 7}{\displaystyle 8}\right)^{k-2},\allowbreak \ k=2,\allowbreak 3,\allowbreak 4,\allowbreak \ldots\)。（II）\(\displaystyle 16\)。


(23)（I）\(\displaystyle \hat\theta=2\overline X-1\)。（II）\(\displaystyle \hat\theta=\displaystyle \min_{1\leq i\leq n}X_i\)。


\clearpage
\end{document}
